%=============================================================================!
% 3.14  WHEN THE ENERGY CHANGES                                               !
%=============================================================================!
\section{When the energy changes}
\label{sec:en-drift}

\Cref{prop:en-energy-atoms} needs two conditions: every force is minus the
gradient of a potential energy, and that potential energy does not change
with time on its own. This section works out what happens to the total
energy of a group of atoms when either fails, and why that makes the
energy the first thing to check in a simulation of atoms isolated from
their surroundings.

\subsection*{A force that is not a gradient}

Write the force on each atom as a conservative part plus the rest,
$\vb{F}_i = -\grad_i U + \vb{F}_i^{\mathrm{nc}}$. The rest adds to the
power that changes the kinetic energy, $\dd K/\dd t = \sum_i\vb{F}_i
\cdot\vb{v}_i$, but not to the rate of change of $U$, which the chain rule
along a path still gives as $\sum_i\grad_i U\cdot\vb{v}_i$. Adding the two
rates,
\begin{align}
   \frac{\dd E}{\dd t}
   &= \sum_{i=1}^{N}\bigl(-\grad_i U + \vb{F}_i^{\mathrm{nc}}\bigr)\cdot
      \vb{v}_i + \sum_{i=1}^{N}\grad_i U\cdot\vb{v}_i \nonumber\\
   &= \sum_{i=1}^{N}\vb{F}_i^{\mathrm{nc}}\cdot\vb{v}_i ,
      \qquad\text{the gradient terms cancelling,}
   \label{eq:en-dEdt}
\end{align}
the rate at which the non-conservative forces do work, as for the box on
the floor in \cref{eq:en-dEdt-1d}. A drag on every atom, $\vb{F}_i^{
\mathrm{nc}} = -\gamma m_i\vb{v}_i$ with the drag rate $\gamma$ of
\cref{sec:ne-drag}, gives, since $\vb{v}_i\cdot\vb{v}_i = v_i^2$ and
$\sum_i m_i v_i^2 = 2K$,
\begin{equation*}
   \frac{\dd E}{\dd t} = -\gamma\sum_{i=1}^{N} m_i v_i^2 = -2\gamma K ,
\end{equation*}
so the energy can only fall. Chapter~13 adds such a drag on purpose, with
random kicks that make up for it, to hold atoms at a chosen temperature.

The swirl shows what a force does that sometimes adds energy and
sometimes removes it.
\Cref{fig:en-drift} repeats the motion of \cref{fig:en-li} with the swirl
of \cref{fig:en-surface}(b) added to the force. The total energy now
wanders, from \SI{0.0185}{\electronvolt} below its starting value to
\SI{0.0117}{\electronvolt} above it within \SI{1.5}{\pico\second}. The
power of the swirl is $\vb{F}_{\mathrm{sw}}\cdot\vb{v} = c\,(x v_y - y
v_x)$, so the atom gains energy while it moves anticlockwise about the
origin and loses it while it moves clockwise: on the positive $x$ axis, for
instance, $y = 0$ and the power is $c\,x v_y$, positive when the atom
moves up, which is anticlockwise there. Its change agrees with the work done by the swirl, the integral
of \cref{eq:en-dEdt} over time, to within \SI{1e-8}{\electronvolt}.

\begin{figure}[htb]
\centering
\includegraphics{ch03_energy/drift}
\caption{The lithium atom of \cref{fig:en-li}, launched in the same way on
the surface alone (teal) and with the swirl of \cref{fig:en-surface}(b)
added to the force (dark red). The change of the total energy $E = K + U$
from its starting value is shown against time, with the work done by the
swirl, $\int_0^t\vb{F}_{\mathrm{sw}}\cdot\vb{v}\,\dd t'$ (dashed), which it follows to within \SI{1e-8}{\electronvolt}.}
\label{fig:en-drift}
\end{figure}

\subsection*{A potential energy that changes with time}

If the potential energy also depends on time, $U(\vb{r}^N, t)$, because
something outside the atoms is being changed while they move, then $U$
changes even when the atoms stand still. A step in time, $\Delta t$, is
then one more leg of the gradient toolbox's step, and the chain rule along
a path gains one term, the partial derivative with respect to time,
\begin{equation*}
   \frac{\dd U}{\dd t} = \sum_{i=1}^{N}\grad_i U\cdot\vb{v}_i
   + \frac{\partial U}{\partial t} ,
\end{equation*}
so that, with conservative forces, $\dd E/\dd t = \partial U/\partial t$.
A hand that moves the far end of a spring does work on the body at the
near end in this way: with the hand at $X(t)$, measured so that the spring
is slack when the body is at $x = X$, the body has the potential energy $U
= \frac12 k\,(x - X(t))^2$, which changes even while the body stands
still. So does a wall pushed in to squeeze a group of atoms, the picture
behind the pistons with which Chapter~14 controls their pressure.

\Needspace{12\baselineskip}
\begin{practice}
The exact motion of an isolated group of atoms, one that exchanges no
energy with its surroundings, keeps $E$ constant. A simulation of such a
group whose total energy drifts is therefore not following the exact
motion, and there are two common causes. The first is the finite time
step with which the equations of motion are advanced, the subject of
Chapter~9. The second is a force that is not exactly $-\grad_i U$ of the
potential energy reported with it, for which \cref{eq:en-dEdt} gives the
drift directly: an error in the code that computes the forces, a
potential energy cut off abruptly at some distance (Chapter~8), or a machine-learned
model, a program fitted to examples of computed forces (Part~V), that
predicts the forces directly instead of differentiating an
energy, so that nothing constrains their curl. Bigi, Langer and Ceriotti
found that models of the last kind can make several types of molecular
dynamics unstable and the convergence of geometry optimisation, the
search for a minimum of $U$, ill-defined, and that their departure from
energy conservation is difficult to monitor\cite{Bigi2025}. Even so,
recording $E$ in a short run of the isolated system is the first check of any new
model or code.
\end{practice}

\begin{keyidea}
For atoms whose forces are $\vb{F}_i = -\grad_i U$, with $U$ fixed in
time, Newton's laws keep $E = K + U$ constant. Any other part of the
force changes $E$ at the rate at which it does work, $\dd E/\dd t =
\sum_i\vb{F}_i^{\mathrm{nc}}\cdot\vb{v}_i$, and a change of $U$ with time
changes $E$ at the rate $\partial U/\partial t$.
\end{keyidea}

\notebook{03}{vary the strength of the swirl and watch the energy wander}

\begin{selfcheck}
In a simulation of an isolated group of atoms, the total energy rises
steadily. Name two possible causes, and say which equation of this
section describes the second.
\end{selfcheck}
